\documentclass[11pt, dina4, amsfont12]{article}
\topmargin=-1.0cm
\usepackage{graphicx}
\usepackage{epstopdf}
\usepackage{showlabe}
\usepackage{amssymb}
\usepackage{amsmath}
\usepackage{amsthm}
\usepackage{color}
\input psfig.sty
\usepackage[margin=1in]{geometry}
\usepackage{hyperref}

% ******************************* Author Macros *******************************

\newcommand{\pl}[2]{\frac{\partial#1}{\partial#2}}
\newcommand{\bu}{{\bf u} }
\newcommand{\p}{\partial}
\newcommand{\og}{\omega}
\newcommand{\Og}{\Omega}
\newcommand{\fl}[2]{\frac{#1}{#2}}
\newcommand{\dt}{\delta}
\newcommand{\nn}{\nonumber}
\newcommand{\ap}{\alpha}
\newcommand{\bt}{\beta}
\newcommand{\tht}{\theta}
\newcommand{\kp}{\kappa}
\newcommand{\Dt}{\Delta}
\renewcommand{\theequation}{\arabic{section}.\arabic{equation}}
\newcommand{\be}{\begin{equation}}
\newcommand{\ee}{\end{equation}}
\newcommand{\ba}{\begin{array}}
\newcommand{\ea}{\end{array}}
\newcommand{\bea}{\begin{eqnarray}}
\newcommand{\eea}{\end{eqnarray}}
\newcommand{\beas}{\begin{eqnarray*}}
\newcommand{\eeas}{\end{eqnarray*}}
\newtheorem{theorem}{Theorem}[section]
\newtheorem{lemma}{Lemma}[section]
\newtheorem{proposition}{Proposition}[section]
\newtheorem{remark}{Remark}[section]
\newtheorem{assumption}{Assumption}[section]
\newcommand{\bx}{{\bf x} }
\newcommand{\gm}{\gamma}
\newcommand{\vep}{\varepsilon}

% *****************************************************************************

\title{Notes on the proof of the remainder}

\author{
Weijie Huang\thanks{School of Mathematics and Statistics, Beijing Jiaotong University, Beijing 100044, People's Republic of China {\tt (wjhuang@bjtu.edu.cn)}} \ and
Xinran Ruan\thanks{School of Mathematical Sciences, Capital Normal University, Beijing, China, 100048, People's Republic of China {\tt (xinran.ruan@cnu.edu.cn)}.}
}

\begin{document}
\date{}
\maketitle
%%============后续 严格的数值分析================
\begin{assumption}[Trait-direction stability estimates]
\label{ass:trait_stability}
On the time interval or for the stationary state under consideration, assume
that the following estimates hold.

\begin{enumerate}
\item[(B1)] The discrete phase satisfies
\[
    \max_k
    \left(
        |D_\theta^-u_k^\varepsilon|
        +
        |D_\theta^+u_k^\varepsilon|
    \right)
    \le L_T .
\]

\item[(B2)] The amplitude satisfies the weighted first-difference estimate
\[
    \Delta\theta
    \sum_{k=1}^K
    \left(
        E_k^\varepsilon+E_{k+1}^\varepsilon
    \right)
    |D_\theta^+W_{j,k}^\varepsilon|
    \le
    M_T m_j^\varepsilon .
\]

\item[(B3)] The numerical Hamiltonian is bounded on bounded slopes,
\[
    \left|
    \widehat H
    \left(
        D_\theta^-u_k^\varepsilon,
        D_\theta^+u_k^\varepsilon
    \right)
    \right|
    \le
    C_H L_T^2 .
\]

\item[(B4)] The conservative trait flux satisfies the weighted interface
bound
\[
    E_k^\varepsilon
    \left|
    \widehat{\mathcal F}_{j,k+1/2}
    \right|
    \le
    C_F L_T
    \left(
        n_{j,k}^\varepsilon+n_{j,k+1}^\varepsilon
    \right),
\]
and the analogous estimate holds for the interface \(k-1/2\).
\end{enumerate}
\end{assumption}

\begin{remark}
For the Godunov Hamiltonian \(\widehat H_G\) and the upwind flux given above,
(B3) and (B4) hold on bounded slope sets. Indeed,
\[
    |\widehat H_G(p^-,p^+)|
    \le
    |p^-|^2+|p^+|^2,
\]
and the upwind choice avoids uncontrolled ratios of neighboring exponential
weights.
\end{remark}

\begin{lemma}[One-sided control of the weighted trait remainder]
\label{lem:one_sided_remainder_bound}
Assume that Assumption~\ref{ass:trait_stability} holds. Then
\begin{equation}\label{eq:one_sided_remainder_bound}
    \mathcal R_j^{\varepsilon,D}
    \le
    C
    \left(
        L_T^2
        +
        \frac{\varepsilon L_T}{\Delta\theta}
        +
        \frac{\varepsilon^2M_T}{\Delta\theta}
    \right)
    m_j^\varepsilon .
\end{equation}
In particular, for fixed \(\Delta\theta>0\) and \(0<\varepsilon\le1\), there
exists a constant \(C_R=C_R(\Delta\theta,L_T,M_T)\), independent of
\(\varepsilon\), such that
\begin{equation}\label{eq:remainder_CR_bound}
    \mathcal R_j^{\varepsilon,D}
    \le
    C_R m_j^\varepsilon .
\end{equation}
\end{lemma}

\begin{proof}
We estimate the four terms in \(R_{j,k}^\varepsilon\). Throughout the proof,
\(C\) denotes a constant independent of \(\varepsilon\).

For the first term, using
\[
    \delta_\theta^2W_{j,k}^\varepsilon
    =
    \frac{
        D_\theta^+W_{j,k}^\varepsilon
        -
        D_\theta^+W_{j,k-1}^\varepsilon
    }{\Delta\theta},
\]
we obtain
\[
\begin{aligned}
&
\varepsilon^2
\Delta\theta
\sum_{k=1}^K
\frac{E_k^\varepsilon}{D_k}
\left|
\delta_\theta^2W_{j,k}^\varepsilon
\right|
\\
&\qquad
\le
C
\frac{\varepsilon^2}{\Delta\theta}
\Delta\theta
\sum_{k=1}^K
\left(
    E_k^\varepsilon+E_{k+1}^\varepsilon
\right)
\left|
D_\theta^+W_{j,k}^\varepsilon
\right|
\\
&\qquad
\le
C
\frac{\varepsilon^2M_T}{\Delta\theta}
m_j^\varepsilon .
\end{aligned}
\]
For the second term, (B1) gives
\[
    |\delta_\theta^2u_k^\varepsilon|
    =
    \left|
    \frac{
        D_\theta^+u_k^\varepsilon
        -
        D_\theta^+u_{k-1}^\varepsilon
    }{\Delta\theta}
    \right|
    \le
    \frac{2L_T}{\Delta\theta}.
\]
Therefore,
\[
\begin{aligned}
&
\varepsilon
\Delta\theta
\sum_{k=1}^K
\frac{n_{j,k}^\varepsilon}{D_k}
\left|
\delta_\theta^2u_k^\varepsilon
\right|
\\
&\qquad
\le
C
\frac{\varepsilon L_T}{\Delta\theta}
m_j^\varepsilon .
\end{aligned}
\]
For the Hamiltonian term, (B3) yields
\[
\begin{aligned}
&
\Delta\theta
\sum_{k=1}^K
\frac{n_{j,k}^\varepsilon}{D_k}
\left|
\widehat H
\left(
    D_\theta^-u_k^\varepsilon,
    D_\theta^+u_k^\varepsilon
\right)
\right|
\\
&\qquad
\le
C L_T^2 m_j^\varepsilon .
\end{aligned}
\]
Finally, for the flux term,
\[
\begin{aligned}
&
2\varepsilon
\Delta\theta
\sum_{k=1}^K
\frac{E_k^\varepsilon}{D_k}
\left|
\mathcal D_\theta
\widehat{\mathcal F}_{j,k}
\right|
\\
&\qquad
\le
C\varepsilon
\sum_{k=1}^K
E_k^\varepsilon
\left(
    |\widehat{\mathcal F}_{j,k+1/2}|
    +
    |\widehat{\mathcal F}_{j,k-1/2}|
\right)
\\
&\qquad
\le
C
\frac{\varepsilon L_T}{\Delta\theta}
m_j^\varepsilon .
\end{aligned}
\]
Combining the four estimates gives \eqref{eq:one_sided_remainder_bound}.
The bound \eqref{eq:remainder_CR_bound} follows for fixed
\(\Delta\theta\) and \(0<\varepsilon\le1\).
\end{proof}


\end{document}